\subsection{Frobenius stability and the Boolean kernel partition}\label{sec:binary:kernel}

Fix $\ell\geq1$ and put $n:=2^\ell$. Let $I_\ell=\{0,1\}^\ell$, and let
$K_b(X)=\prod_{i=0}^{\ell-1}(X^{2^i}+b_i)$ denote the Boolean kernel from Definition~\ref{def:intro:kernel}.

\begin{lemma}[Binary coefficients of a kernel]\label{lem:binary:binary-coeffs}
For every $b\in I_\ell$, the polynomial $K_b$ has all coefficients in
the prime subfield $\F_2\subseteq\F$.
\end{lemma}
\begin{proof}
Every factor $X^{2^i}+b_i$ has coefficients in $\F_2$ because
$b_i\in\{0,1\}$. A finite product of polynomials with coefficients
in a subfield again has coefficients in that subfield, by the
finite-sum formula for multiplication coefficients.
\end{proof}

\begin{lemma}[Kernel coefficient Frobenius invariance]\label{lem:binary:kernel-coeff-frob}
For $b\in I_\ell$ and $j\geq0$,
\[
 K_b^{\sigma_j}(X)=K_b(X).
\]
\end{lemma}
\begin{proof}
By Lemma~\ref{lem:binary:binary-coeffs}, each coefficient is
$0$ or $1$. Both are fixed by $\sigma_j$, so the coefficientwise
Frobenius operation leaves the polynomial unchanged.
\end{proof}

\begin{theorem}[Kernel Frobenius substitution]\label{thm:binary:kernel-square}
For every $b\in I_\ell$ and $j\geq0$,
\[
 \boxed{K_b(X)^{2^j}=K_b(X^{2^j}).}
\]
\end{theorem}
\begin{proof}
Lemma~\ref{lem:binary:poly-frob} gives
$K_b(X)^{2^j}=K_b^{\sigma_j}(X^{2^j})$.
By Lemma~\ref{lem:binary:kernel-coeff-frob} the latter equals
$K_b(X^{2^j})$.
\end{proof}

\begin{definition}[Kernel coefficient Frobenius]\label{def:binary:table-frob}
For a coefficient table $c:I_\ell\to\F$ and $j\geq0$, define
$c^{\sigma_j}(b):=c(b)^{2^j}$ and
$U_c(X):=\sum_{b\in I_\ell}c(b)K_b(X)$.
\end{definition}

\begin{theorem}[Frobenius covariance of kernel encoding]\label{thm:binary:encoding-frob}
For every coefficient table $c$ and $j\geq0$,
\[
 \boxed{U_c(X)^{2^j}=U_{c^{\sigma_j}}(X^{2^j}).}
\]
\end{theorem}
\begin{proof}
Repeatedly apply Lemma~\ref{lem:binary:square} to the finite sum
defining $U_c$:
\[
 U_c(X)^{2^j}
 =\sum_b c(b)^{2^j}K_b(X)^{2^j}.
\]
Theorem~\ref{thm:binary:kernel-square} turns each kernel factor
into $K_b(X^{2^j})$, giving
$\sum_b c^{\sigma_j}(b)K_b(X^{2^j})$,
which is $U_{c^{\sigma_j}}(X^{2^j})$.
\end{proof}

\begin{lemma}[Kernel partition sum]\label{lem:binary:partition}
\[
 \boxed{\sum_{b\in I_\ell}K_b(X)=1.}
\]
\end{lemma}
\begin{proof}
Distribute the sum over all bit strings:
\[
 \sum_{b\in I_\ell}\prod_{i=0}^{\ell-1}(X^{2^i}+b_i)
 =\prod_{i=0}^{\ell-1}
    \bigl((X^{2^i}+0)+(X^{2^i}+1)\bigr).
\]
Inside each parenthesis, $X^{2^i}+X^{2^i}=0$ in
characteristic two, leaving the scalar $1$. The product of
$\ell$ copies of $1$ is $1$.
\end{proof}

\begin{corollary}[All-one kernel coordinates of the constant polynomial]\label{cor:binary:one-coeffs}
Let $\mathbf1:I_\ell\to\F$ be the constant table with value $1$.
Then $U_{\mathbf1}(X)=1$. Moreover the coordinates of $1$
in the kernel basis are uniquely all equal to $1$.
\end{corollary}
\begin{proof}
By definition $U_{\mathbf1}=\sum_bK_b=1$,
using Lemma~\ref{lem:binary:partition}. Uniqueness follows
from the kernel basis theorem, Proposition~\ref{prop:intro:basis}.
\end{proof}

\begin{lemma}[Two complementary lowest-bit kernels]\label{lem:binary:complementary}
For every $b'\in I_{\ell-1}$,
\[
 K_{(0,b')}(X)+K_{(1,b')}(X)=K_{b'}^{(\ell-1)}(X^2).
\]
\end{lemma}
\begin{proof}
By Lemma~\ref{lem:intro:recursion},
$K_{(0,b')}=XK_{b'}^{(\ell-1)}(X^2)$ and
$K_{(1,b')}=(X+1)K_{b'}^{(\ell-1)}(X^2)$.
Add and factor out the common polynomial. The scalar
$X+(X+1)$ equals $1$ because $X+X=0$.
\end{proof}

\begin{remark}[One important characteristic restriction]
Lemma~\ref{lem:binary:partition} is false over a general
odd-characteristic field. The exact all-characteristic identity
is Proposition~\ref{prop:kbalgebra:partition}. We will use
$\sum_bK_b=1$ to identify the degeneracy of an \emph{unanchored}
residual polynomial, so the characteristic assumption cannot
be dropped.
\end{remark}
